\documentclass[12pt]{report}

% =========================
% APA 7 Formatting
% =========================
\usepackage[american]{babel}
\usepackage{csquotes}
\usepackage[style=apa, backend=biber]{biblatex}
\DeclareLanguageMapping{american}{american-apa}
\addbibresource{references.bib}

\usepackage{setspace}
\onehalfspacing

\usepackage[margin=1in]{geometry}
\usepackage{indentfirst}

\usepackage{titlesec}
\titleformat{\chapter}{\normalfont\bfseries\centering}{\thechapter}{1em}{}

% =========================
% Mathematics Packages
% =========================
\usepackage{amsmath, amssymb, amsthm, mathtools}
\usepackage{bm}
\usepackage{physics}
\usepackage{siunitx}

\numberwithin{equation}{chapter}

% =========================
% Theorem Environments
% =========================
\theoremstyle{definition}
\newtheorem{definition}{Definition}[chapter]

\theoremstyle{plain}
\newtheorem{theorem}[definition]{Theorem}
\newtheorem{lemma}[definition]{Lemma}
\newtheorem{corollary}[definition]{Corollary}
\newtheorem{proposition}[definition]{Proposition}

\theoremstyle{remark}
\newtheorem{remark}[definition]{Remark}
\newtheorem{example}[definition]{Example}

% =========================
% Hyperlinks
% =========================
\usepackage[colorlinks=true, linkcolor=blue, citecolor=blue]{hyperref}

% =========================
% Title Page
% =========================
\title{Advanced Mathematical and Financial Models}
\author{Joseph Burton}
\date{\today}

\begin{document}

\maketitle
\tableofcontents
\newpage

% =====================================================
\chapter{Foundations of Mathematical Analysis}

\section{Managerial Accounting Vs Financial Accounting}

Financial Accounting deals with outside constituents: Investors, Lenders and
Government. Managerial Accounting exist for the insiders Managers and
Employees. The former are making choices \textbf{about} the company and the
latter are making choices \textbf{for} the company. FA is about the past and
the whole company. MA is mostly about the future and deals with the local and
not th global!

\begin{definition}[Operation Cash Flow]
    Operating Cash Flow (OCF) measures the cash generated by a company's core business activities,
    excluding long-term capital investments or financing costs.

    \[ OCF = NI + NCE = IWC\]

    \textbf{NI} = Net Income \\
    \textbf{NCE} = Non-Cash Expenses \\
    \textbf{IWC} = Increases in Working Capital

    \[ OCF = EBIT + Depreciation - Taxes \]

    \textbf{EBIT} = Earnings Before Intrest and Taxes
\end{definition}

\begin{definition}[Capital Expendenture]
    Capital expenditure (CapEx) represents funds used by a company to acquire, upgrade, and maintain physical assets
    like property, buildings, or equipment to increase operational capacity or longevity.

    \[ CapEx = \Delta PP\&E + Current Depreciation  \]
\end{definition}

\begin{definition}[Change in Net Working Capital]
    Change in Net Working Capital (NWC) is the difference between a company's
    NWC at the end of a period and its NWC at the beginning of that period

    \[ NWC_{end} - NWC_{start} \]

    It measures how operational liquidity changes, with a positive change signaling
    expansion and a negative change often indicating increased cash flow from
    inventory or receivable reductions

    \[ NWC = (Current Assets - Current Liabilities) \]
\end{definition}

% =====================================================
\chapter*{References}
\addcontentsline{toc}{chapter}{References}
\printbibliography

\end{document}